\begin{table}[t]
\centering\small
\caption{Estudio de ablación. Todas las corridas usan el mismo presupuesto (3000 pasos, lote 16, recortes de 256\,px). La selección se hace con el IoU de agua en \emph{validación} (el mejor, en negrita); las columnas de prueba y Bolivia se reportan solo como referencia.}
\label{tab:ablacion}
\setlength{\tabcolsep}{4pt}
\resizebox{\linewidth}{!}{\begin{tabular}{lllcrcccc}
\toprule
\textbf{Corrida} & \textbf{Datos} & \textbf{Pérdida} & \textbf{Arquitectura} & \textbf{Parám. (M)} & \textbf{IoU val} & \textbf{IoU test} & \textbf{F1 test} & \textbf{IoU Bolivia} \\
\midrule
E1 & dorado & BCE & con saltos, prof. 4 & 7.76 & 0.615 & 0.640 & 0.781 & 0.408 \\
E2 & dorado + débil & BCE & con saltos, prof. 4 & 7.76 & 0.598 & 0.641 & 0.781 & 0.716 \\
E3 & dorado & Dice & con saltos, prof. 4 & 7.76 & 0.630 & 0.661 & 0.796 & 0.615 \\
E4 & dorado & BCE + Dice & con saltos, prof. 4 & 7.76 & 0.621 & 0.656 & 0.792 & 0.518 \\
E6 & dorado & BCE con peso en bordes & con saltos, prof. 4 & 7.76 & 0.624 & 0.665 & 0.799 & 0.598 \\
E5a & dorado & Dice & sin saltos, prof. 4 & 6.98 & 0.634 & 0.649 & 0.787 & 0.594 \\
E5b & dorado & Dice & con saltos, prof. 3 & 1.93 & \textbf{0.641} & 0.663 & 0.798 & 0.572 \\
\bottomrule
\end{tabular}}
\end{table}
